\begin{center}\begin{tikzpicture}[font=\small,>=stealth,line join=round]
\draw[gray!40,rounded corners=3pt] (0,0) rectangle (14.2,8.3);\node at (7.1,7.700000000000001) {同一高度的两个中心，半转两次得到水平平移};
\begin{scope}[shift={(2.463636363636364,3.9499999999999997)},x=0.5454545454545454cm,y=1.5333333333333332cm]\draw[->,gray] (-2.5,0)--(8.5,0) node[right]{$x$};\draw[->,gray] (0,-1.5)--(0,1.5) node[above]{$y$};\draw[blue,thick,solid,domain=-2:8,samples=180] plot(\x,{sin(deg(pi*\x/3))});\fill[orange!85!black] (1,0.8660254037844386) circle(2.2pt);\fill[orange!85!black] (-1,-0.8660254037844386) circle(2.2pt);\fill[orange!85!black] (7,0.8660254037844386) circle(2.2pt);\end{scope}\node[anchor=west,align=left,text width=5.0cm] at (8.4,4.0) {图取 $f=\sin\frac{\pi x}{3}$\\[9pt]中心 $(0,0)$、$(3,0)$\\[9pt]$1\to-1\to7$\\[9pt]高度依次 $\frac{\sqrt3}2,-\frac{\sqrt3}2,\frac{\sqrt3}2$\\[9pt]相差 $6$，回到同高};
\node[text width=12.6cm,align=center] at (7.1,.65) {正文对任意函数的配对关系推出 $f(x+6)=f(x)$；6只保证是一个周期。};
\end{tikzpicture}\end{center}
